\begin{tikzpicture}[
  font=\scriptsize,
  kotak/.style={draw, rounded corners, minimum height=1.0cm, text width=1.45cm,
                align=center, font=\scriptsize\bfseries},
  tanya/.style={text width=1.6cm, align=center, font=\tiny},
  panah/.style={-{Stealth[length=2mm]}, thick}
]
  \node[kotak, fill=praditagreen!15] (daftar) at (0,0)
    {Paket kerja\\{\tiny\mdseries dari Bab 11}};
  \node[kotak, fill=blue!8] (dep) at (2.05,0)
    {Dependensi\\{\tiny\mdseries syarat mutlak}};
  \node[kotak, fill=blue!8] (nilai) at (4.1,0)
    {Nilai dan risiko\\{\tiny\mdseries skor}};
  \node[kotak, fill=orange!12] (kap) at (6.15,0)
    {Kapasitas\\{\tiny\mdseries batas per tahun}};
  \node[kotak, fill=blue!8] (angg) at (8.2,0)
    {Anggaran\\{\tiny\mdseries titik keputusan}};
  \node[kotak, fill=green!10] (peta) at (10.25,0)
    {Peta jalan\\{\tiny\mdseries 2026 sampai 2030}};

  \draw[panah] (daftar) -- (dep);
  \draw[panah] (dep) -- (nilai);
  \draw[panah] (nilai) -- (kap);
  \draw[panah] (kap) -- (angg);
  \draw[panah] (angg) -- (peta);

  \node[tanya, anchor=north] at (0,-0.7) {6 sampai 10 paket beserta matriks dependensi};
  \node[tanya, anchor=north] at (2.05,-0.7) {Apa yang harus selesai lebih dahulu?};
  \node[tanya, anchor=north] at (4.1,-0.7)
    {Mana yang paling bermanfaat dan paling menurunkan risiko?};
  \node[tanya, anchor=north] at (6.15,-0.7)
    {Berapa proyek besar yang sanggup berjalan dalam setahun?};
  \node[tanya, anchor=north] at (8.2,-0.7) {Kapan dana tahun berikutnya diputuskan?};
  \node[tanya, anchor=north] at (10.25,-0.7) {Urutan, transisi, dan titik keputusan};
\end{tikzpicture}
